% ============================================================
%  ELSEVIER — Computers & Industrial Engineering
%  Full journal template — single self-contained main.tex
%  Replicates: Elsevier two-column layout, fonts, colours,
%  header/footer, section style, article-info block, abstract
%  box with keywords, numbered equations, booktabs tables,
%  nomenclature table, large wide result tables, TikZ figures
%  (line chart with error bars, supply-chain schematic),
%  numbered citation style (Vancouver/Elsevier), references.
% ============================================================

% ── Engine ──────────────────────────────────────────────────
\documentclass[10pt,twocolumn]{article}

% ── Core packages ───────────────────────────────────────────
\usepackage[T1]{fontenc}
\usepackage[utf8]{inputenc}

% ── Fonts ────────────────────────────────────────────────────
% Elsevier body: Times New Roman equivalent
\usepackage{mathptmx}
\usepackage[scaled=0.90]{helvet}
\usepackage{courier}
\renewcommand{\sfdefault}{phv}

% ── Page geometry ────────────────────────────────────────────
% CIE: ~170 mm text width, 247 mm text height, two columns
\usepackage[
  paperwidth  = 210mm,
  paperheight = 297mm,
  top         = 22mm,
  bottom      = 25mm,
  left        = 19mm,
  right       = 19mm,
  columnsep   = 5.5mm,
  headsep     = 7pt,
  footskip    = 14pt
]{geometry}

% ── Colours ─────────────────────────────────────────────────
\usepackage[dvipsnames,table]{xcolor}
% Elsevier blue used in title, section rules, hyperlinks
\definecolor{elsblue}{RGB}{0,84,166}
% Header/footer text grey
\definecolor{elsgrey}{RGB}{90,90,90}
% Light rule colour
\definecolor{elsrule}{RGB}{160,160,160}
% Highlight colour for section heading rule (thin red accent)
\definecolor{elsred}{RGB}{200,0,0}

% ── Header / footer ─────────────────────────────────────────
\usepackage{fancyhdr}
\pagestyle{fancy}
\fancyhf{}
\fancyhead[L]{\footnotesize\itshape\color{elsgrey}
  M.R.G.~Samani et al.}
\fancyhead[R]{\footnotesize\itshape\color{elsgrey}
  \textit{Computers \& Industrial Engineering}~130~(2019)~450--471}
\fancyfoot[C]{\small\thepage}
\renewcommand{\headrulewidth}{0.5pt}
\renewcommand{\headrule}{{\color{elsrule}%
  \hrule width\headwidth height\headrulewidth}}

\fancypagestyle{firstpage}{
  \fancyhf{}
  \fancyhead[L]{%
    \footnotesize\color{elsgrey}
    \textit{Computers \& Industrial Engineering}~130~(2019)~450--471}
  \fancyfoot[C]{\small\thepage}
  \renewcommand{\headrulewidth}{0.5pt}
  \renewcommand{\headrule}{{\color{elsrule}%
    \hrule width\headwidth height\headrulewidth}}
}

% ── Line spacing & paragraph ─────────────────────────────────
\usepackage{setspace}
\setstretch{1.04}
\setlength{\parindent}{1em}
\setlength{\parskip}{2pt plus 1pt}

% ── Section heading style ────────────────────────────────────
\usepackage{titlesec}
% Level-1: bold, 10.5 pt, left-flush, thin rule below
\titleformat{\section}
  {\normalfont\fontsize{10.5}{13}\bfseries}
  {\thesection.}
  {0.5em}{}
  [\vspace{1pt}\color{elsrule}\rule{\columnwidth}{0.4pt}]
\titlespacing{\section}{0pt}{9pt plus 2pt}{4pt}

% Level-2: bold italic, 10 pt
\titleformat{\subsection}
  {\normalfont\fontsize{10}{12}\bfseries\itshape}
  {\thesubsection.}
  {0.5em}{}
\titlespacing{\subsection}{0pt}{7pt plus 2pt}{3pt}

% Level-3: italic run-in
\titleformat{\subsubsection}[runin]
  {\normalfont\fontsize{10}{12}\itshape}
  {\thesubsubsection.}
  {0.5em}{}[.\quad]
\titlespacing{\subsubsection}{0pt}{5pt}{0pt}

% ── Math ─────────────────────────────────────────────────────
\usepackage{amsmath,amssymb,amsthm}

% ── Algorithm ────────────────────────────────────────────────
\usepackage{algorithm}
\usepackage{algorithmic}

% ── Figures ──────────────────────────────────────────────────
\usepackage{graphicx,float}
\usepackage{tikz}
\usepackage{pgfplots}
\pgfplotsset{compat=1.18}
\usetikzlibrary{arrows.meta,positioning,shapes.geometric,
                decorations.pathmorphing,patterns,calc,
                shapes.arrows,fit,backgrounds}

% ── Tables ───────────────────────────────────────────────────
\usepackage{booktabs,array,tabularx,multirow,longtable,makecell}
\newcommand{\tabfont}{\fontsize{8}{9.5}\selectfont}
\newcommand{\headfont}{\fontsize{8.5}{10}\selectfont\bfseries}
\renewcommand{\arraystretch}{1.1}

% ── Caption style ────────────────────────────────────────────
\usepackage[
  format       = plain,
  labelfont    = {bf, small},
  textfont     = {small},
  labelsep     = period,
  justification= raggedright,
  singlelinecheck = false
]{caption}

% ── Lists ────────────────────────────────────────────────────
\usepackage{enumitem}
\setlist[itemize]{leftmargin=*,topsep=2pt,itemsep=1pt,parsep=0pt,
                  label=\textbullet}
\setlist[enumerate]{leftmargin=*,topsep=2pt,itemsep=1pt,parsep=0pt}

% ── Hyperlinks & numbered citations ─────────────────────────
\usepackage[
  colorlinks = true,
  linkcolor  = elsblue,
  citecolor  = elsblue,
  urlcolor   = elsblue
]{hyperref}
\usepackage{url}
% Elsevier uses numbered, sorted citations (Vancouver style)
\usepackage[numbers,sort&compress]{natbib}
\setlength{\bibsep}{2pt plus 1pt}

% ── Miscellaneous ────────────────────────────────────────────
\usepackage{microtype}
\usepackage{balance}

% ── Utility: double-rule for article-info top bar ───────────
\newcommand{\elstopbar}{%
  \noindent\color{elsblue}\rule{\linewidth}{1.5pt}\\[-3pt]
  \color{elsblue}\rule{\linewidth}{0.5pt}\vspace{4pt}}
\newcommand{\elsbotbar}{%
  \vspace{3pt}\noindent\color{elsblue}%
  \rule{\linewidth}{0.5pt}\\[-3pt]
  \color{elsblue}\rule{\linewidth}{1.5pt}}

% ============================================================
%  DOCUMENT
% ============================================================
\begin{document}

% ── Full-width title block ───────────────────────────────────
\twocolumn[{%

% -------- journal strip ------------------------------------
\noindent
\begin{minipage}[t]{22mm}
  \includegraphics[height=18mm]{example-image}% Elsevier logo
\end{minipage}
\hfill
\begin{minipage}[t]{102mm}
  \centering
  \vspace{2pt}
  {\footnotesize Contents lists available at
   \textcolor{elsblue}{ScienceDirect}\par}
  \vspace{3pt}
  {\fontsize{16}{20}\selectfont\sffamily
   \textbf{Computers \& Industrial Engineering}\par}
  \vspace{2pt}
  {\footnotesize journal homepage:~%
   \textcolor{elsblue}{\url{www.elsevier.com/locate/caie}}\par}
\end{minipage}
\hfill
\begin{minipage}[t]{22mm}
  \includegraphics[height=18mm]{example-image}% CIE cover thumbnail
\end{minipage}

\vspace{8pt}
\noindent\color{elsrule}\rule{\linewidth}{0.5pt}
\vspace{8pt}

% -------- article title ------------------------------------
{\fontsize{17}{21}\selectfont
A [multilateral/multi-objective] perspective towards [supply chain
network design] in an uncertain environment: Methodology and
implementation\par}

\vspace{6pt}

% -------- authors ------------------------------------------
{\fontsize{10.5}{13}\selectfont
First Author\textsuperscript{a},
Second Author\textsuperscript{a,}\textsuperscript{*},
Third Author\textsuperscript{a}\par}

\vspace{3pt}
{\footnotesize\itshape
\textsuperscript{a}School of Industrial Engineering, University Name,
University Ave, City, Country\par}

\vspace{6pt}

% -------- article info two-panel block ----------------------
\noindent
\begin{minipage}[t]{48mm}
  \fontsize{8.5}{11}\selectfont
  \textbf{ARTICLE INFO}\\[4pt]
  \textit{Keywords:}\\
  [Keyword 1]\\
  [Keyword 2]\\
  [Keyword 3]\\
  [Keyword 4]
\end{minipage}%
\hfill
\begin{minipage}[t]{114mm}
  \fontsize{8.5}{11}\selectfont
  \textbf{ABSTRACT}\\[4pt]
  Motivated by a real-life case study, this paper proposes a
  [multilateral perspective / novel framework] for [supply chain
  network / domain] design as an intricate decision-making problem.
  The proposed approach utilizes a novel [multi-objective mathematical
  model] by incorporating both quantitative and qualitative factors to
  comprehensively model the investigated case study. In the proposed
  model, the quantitative factors aim to [minimize / maximize]
  [objective 1] and [objective 2]. Moreover, the number of [facilities],
  which are capable of [performing function], are [maximized / minimized]
  through a formulation based upon qualitative criteria. Moreover, to
  cope with the uncertainty, a robust optimization approach is utilized,
  and an interactive fuzzy solution approach is applied to solve the
  proposed multi-objective model. The computational results demonstrate
  the capability of the proposed model in designing the [network] in
  practice. Moreover, some sensitivity analyses are conducted to show
  the findings observed in a real case study. The results show that the
  decision makers should be aware of [quality-based / cost-based]
  requirements, because of significant impact on [facility location /
  resource allocation], besides the [cost / quality]-oriented
  consideration. Finally, based on the finding, helpful managerial
  insights are provided for the decision makers in the [domain] network
  design.
\end{minipage}

\vspace{6pt}
\noindent\color{elsrule}\rule{\linewidth}{0.5pt}
\vspace{6pt}

% -------- received / accepted dates -------------------------
{\footnotesize
Received [DD Month YYYY];
Received in revised form [DD Month YYYY];
Accepted [DD Month YYYY]\\
Available online [DD Month YYYY]\\
[ISSN] / \copyright\ 202X [Publisher]. All rights reserved.\par}

\vspace{10pt}
}]% end \twocolumn[...]

\thispagestyle{firstpage}

% ============================================================
\section{Introduction}
% ============================================================

The world fast-growing population has resulted in an increase in
[diseases / costs / demand] which redouble the importance of
[domain] considerations. Among the items used, [product] and its
sub-products are crucial commodities in any [sector] because they
cannot be replaced by any alternatives, and each of them has its
special usage. In addition, [domain]-related cost owns a significant
part of total [sector] expenditures. By 2040, [relevant projection]
expenditure is expected to have an approximate growth of [X]\%
\cite{Author2015}. Therefore, [domain] supply chain (DSC)
management due to its inherent importance in [sector] has become an
important decision-making problem and drawn attention from
researchers.

[Domain] supply chains are complex and difficult to manage owing to
the fact that both supply and demand are tainted with uncertainty.
Since [source] is the only source of [product] with almost
unpredictable behaviors, [product] supply is an unavoidably uncertain
parameter. According to [Organization], on average, only [X]\% of
potential [suppliers] supplied [product] in [year] \cite{WHO2016}.

Generally, uncertain environments can be defined in three types
namely stochastic, fuzzy and robust each of which are dealt with
their corresponding modeling approaches \cite{Tofighi2016}. It is
noteworthy that the stochastic programming approach comes to the
application when the adequate information about distribution
functions of random data is available. However, the robust approach
is applicable in the situation of little knowledge about the historical
data \cite{Torabi2015}. Consequently, it is of particular importance
to make robust decisions, especially for issues with a long-term
horizon.

In this relation, this study sets out to give proper responses to the
following main questions:

\begin{enumerate}
  \item What is the influential qualitative and quantitative factors
        that affect the decisions in [domain] network design? Which
        approach is appropriate to cope with? And how can make a
        tradeoff these factors to effective and efficient performance
        for [domain]?
  \item Which aspects of [domain] are tainted by uncertainty? How to
        formulate the uncertainty in the proposed model? And what are
        the impacts of uncertainty on the proposed model?
  \item How to embed specific characteristics of [product] like a
        shelf life time, different [policies], and optimal inventory
        level of [product] in the proposed model?
\end{enumerate}

In this paper, we utilize effective multi-attribute group
decision-making (MAGDM) approach to account for a variety of
qualitative attributes like [attribute 1], [attribute 2], and
[attribute 3] which are evaluated under the preferences of a group of
beneficiaries as the decision makers (DMs). A three-phase
multi-attribute multi-objective robust (3P-MAMOR) approach is also
devised to cope with the uncertain environment, and then, an
interactive fuzzy approach is provided to handle the objectives
multiplicity \cite{Samani2018,Torabi2008}.

The rest of the paper is organized as follows. Section~2 reviews the
related literature. In Section~3, the solution methodology is
described, and the problem description and mathematical formulation
are provided in Section~4. Section~5 describes the case study in more
details, and Section~6 outlines the computational experiments.
Finally, Section~7 provides discussions, concluding remarks and future
research directions.

% ============================================================
\section{Literature Review}
% ============================================================

In this section, an extant body of literature is investigated in the
form of a structured review. According to the review by
\citet{Belien2012}, the studies were conducted based on the issues
defined on hierarchical levels. Another survey was provided by
\citet{Osorio2015} classifying the papers based on the echelons of
the supply chain. According to Table~\ref{tab:nomenclature}, the
papers are categorized based on hierarchical level, planning horizon,
mathematical modelling approach, time periods, solution technique,
type of [product], objective function multiplicity, performance
measure, and uncertainty modelling approach.

\subsection{[Stage 1] in [Domain] Supply Chains}

Some papers have considered the issues on [stage 1] level
\cite{Author2015a,Author2012,Author2009}. Different [policies] were
evaluated by a simulation model developed by \citet{Author2010}.
Results indicated that [collecting / producing] as much as possible
was not necessarily convenient. \citet{Author2017} addressed [domain]
network design while considering social aspects.

\subsection{[Stage 2] in [Domain] Supply Chains}

A number of works have accounted for [stage 2] planning for [product]
\cite{Author2010b,Author2007}. A decision support system (DSS) for
[product] production planning in regional centers was addressed by
\citet{Author2010b} through an integer non-linear programming model.

\subsection{Uncertainty in [Domain] Supply Chains}

To the best of our knowledge, the majority of the studies which have
addressed the relevant issues in an uncertain environment and to
tackle uncertainty, stochastic programming approach is widely applied
\cite{Author2014,Author2009,Author2017b}. Among them, robust
optimization approaches and its derivatives have got scant attention
\cite{Author2014b,Author2017c,Author2017d}. The aforementioned
literature indicates that the inclusion of both quantitative and
qualitative attributes while designing a [domain] network robustly
performing in an uncertain environment is a critical factor which has
not got apposite attention. To fill this gap, this paper proposes a
realistic approach to investigate the problem from 3P-MAMOR
strategic decision-making perspective which is in line with the
development of previous studies in the literature by addressing the
following main contributions:

\begin{itemize}
  \item A novel multi-objective model for the simultaneous
        consideration of [objective 1] (by investigating quantitative
        attributes), [objective 2], and most favorable locations for
        [facilities] (based on qualitative factors evaluated by MAGDM
        approach);
  \item A realistic multi-product multi-period [domain] network by
        considering three key issues (the design and location of
        facilities, perishability of [products], and uncertainty in
        critical parameters);
  \item A three-phase multi-attribute multi-objective robust approach
        to cope with a group decision-making network design problem
        in an uncertain environment;
  \item Using a real case study of [Location] to illustrate the
        applicability of the proposed model and solution.
\end{itemize}

% ============================================================
\section{Methodology}
% ============================================================

In this section, we suggest a 3P-MAMOR approach to appropriately
design the [domain] supply chain network. The proposed approach
accounts for all information about attributes and qualitative factors
involving in facility location decisions by an effective MAGDM
technique. It also protects the proposed multi-objective model
against uncertainties by applying a robust optimization approach and
devises an interactive fuzzy approach, called TH
\cite{Samani2018,Torabi2008}, to solve the model.

The schematic of the proposed 3P-MAMOR framework can be
demonstrated in Fig.~\ref{fig:framework}. In phase~1
(identification), most appropriate locations for [facilities] among
the candidate ones are identified. In phase~2 (protection), we
confront an unavoidable uncertainty of the parameters and protect the
solutions. In phase~3 (integration), the cost-oriented and
quality-oriented attitudes are integrated into a single objective
structure.

\subsection{Phase~1: The Multi-Attribute Group Decision-Making
Approach}

Decisions on the locations of facilities must be made with respect to
some critical factors which can be [type] considerations. Indeed, we
have to deal with MADM environment in which each qualitative
attribute is evaluated by employing MAGDM perspective based on the
preferences of a group of DMs, each of which differentiates in giving
weights to the attributes.

\subsubsection{Method Framework}
The value of alternative $A_i$ on criterion $C_j$ is shown by the
interval-valued intuitionistic fuzzy setting (IVIFS) where
$A_i = \{\langle C_j, [\mu_{A_iL}(C_j), \mu_{A_iU}(C_j)],
[\vartheta_{A_iL}(C_j), \vartheta_{A_iU}(C_j)]\rangle \mid C_j \in C\}$,
such that $0 \le \mu_{A_iU}(C_j) + \vartheta_{A_iU}(C_j) \le 1$,
$\mu_{A_iL}(C_j) \ge 0$, $\vartheta_{A_iL}(C_j) \ge 0$,
$i = 1,2,\ldots,m$ and $j = 1,2,\ldots,n$. The IVIFS is defined by
$r_{ij} = \langle [a_{ij}, b_{ij}], [c_{ij}, d_{ij}]\rangle$
which represents a pair of intervals for each $C_j \in C$. Thus,
$R^k$ implies on the interval-valued intuitionistic fuzzy decision
matrix $k$ ($k = 1, 2, \ldots, t$) as follows \cite{Ye2013}:

\begin{equation}
R^k =
\begin{pmatrix}
\langle [a^k_{11}, b^k_{11}], [c^k_{11}, d^k_{11}]\rangle
  & \cdots
  & \langle [a^k_{1n}, b^k_{1n}], [c^k_{1n}, d^k_{1n}]\rangle \\
\vdots & \ddots & \vdots \\
\langle [a^k_{m1}, b^k_{m1}], [c^k_{m1}, d^k_{m1}]\rangle
  & \cdots
  & \langle [a^k_{mn}, b^k_{mn}], [c^k_{mn}, d^k_{mn}]\rangle
\end{pmatrix}
\label{eq:ivifs}
\end{equation}

\subsubsection{The Weights of DMs}
Based on matrix $R^k$ introduced above, the score matrix as well as
the average score matrix for each DM can also be calculated
\cite{Ye2013}; and the collective correlation coefficient can be
defined as:

\begin{equation}
E_k = \sum_{i=1}^{m}
  \frac{\sum_{j=1}^{n} S^k_{ij} S^*_{ij}}
  {\sqrt{\sum_{j=1}^{n}(S^k_{ij})^2}
   \sqrt{\sum_{j=1}^{n}(S^*_{ij})^2}}
\label{eq:ek}
\end{equation}

where $S^k_{ij}$ and $S^*_{ij}$ represent the score matrix and the
average score matrix, respectively. Finally, the weight vector
$\Gamma = (\Gamma_1, \Gamma_2, \ldots, \Gamma_t)$ for DMs is
constructed by:

\begin{equation}
\Gamma_k = \frac{E_k}{\sum_{k=1}^{t} E_k}
\label{eq:gamma}
\end{equation}

where $0 \le \Gamma_k \le 1$, $\sum_{k=1}^{t} \Gamma_k = 1$.

\subsubsection{Weights of Attributes and Ranking the Alternatives}
The interval-valued intuitionistic fuzzy decision-making problems can
be developed by means of integrated programming models \cite{Ye2013}
as follows:

\begin{equation}
\begin{aligned}
  \text{Max}\; Q(\omega) &= \sum_{i=1}^{m} \sum_{j=1}^{n} \omega_j S_{ij}\\
  \text{s.t.} \quad &\sum_{j=1}^{n} \omega_j = 1,\quad
                     \omega \in I,\quad \omega_j \ge 0
\end{aligned}
\label{eq:maxq}
\end{equation}

Finally, the alternatives are analyzed using the simple additive
weighting (SAW) method:

\begin{equation}
u(A_i) = \sum_{j=1}^{n} \omega_j S_{ij}
\label{eq:saw}
\end{equation}

where $S_{ij}$ denotes the collective score matrix and $u(A_i)$
denotes the utility of $i$th alternative.

\subsection{Phase~2: Robust Optimization Approach}

In this section, a robust optimization approach is described based
on \citet{Bertsimas2003,Bertsimas2004}. We consider the following
MIP model assuming $X$ is a convex polyhedron:

\begin{equation}
\begin{aligned}
  \text{Min} &\sum_j c_j x_j \\
  \text{s.t.} &\sum_j \tilde{a}_{ij} x_j \le b_i \quad \forall\, i \\
              &x \in X
\end{aligned}
\label{eq:robust_base}
\end{equation}

where $c$ represents the $n$-vector of objective coefficients, and
$b_i$ is the $m$-vector of RHS related to constraint $i$. The
uncertain coefficient $\tilde{a}_{ij}$ takes values in a symmetric
and bounded interval represented as $[a_{ij} - \hat{a}_{ij},\;
a_{ij} + \hat{a}_{ij}]$ in which $\hat{a}_{ij}$ denotes the worst
case. By defining the protection function $\beta(x, \Gamma_i)$ as
Eq.~\eqref{eq:beta}, each uncertain constrain $i$ is rewritten as
Eq.~\eqref{eq:robust_constr}.

\begin{equation}
\beta(x, \Gamma_i) = \max_{\tilde{a}_{ij} \in J_i}
  \left\{\sum_j \hat{a}_{ij} x_j\right\}
  \quad \forall\, i
\label{eq:beta}
\end{equation}

\begin{equation}
\sum_j a_{ij} x_j + \beta(x, \Gamma_i) \le b_i \quad \forall\, i
\label{eq:robust_constr}
\end{equation}

Additionally, the robust model can be formulated using LP duality
as below:

\begin{equation}
\begin{aligned}
  \text{Min} &\sum_j c_j x_j \\
  \text{s.t.} &\sum_j a_{ij} x_j + Z_i \Gamma_i
               + \sum_j p_{ij} \le b_i \quad \forall\, i \\
              &Z_i + p_{ij} \ge \hat{a}_{ij} x_j \quad \forall\, i,j \\
              &p_{ij} \ge 0 \quad \forall\, i,j \\
              &Z_i \ge 0 \quad \forall\, i \\
              &x_j \ge 0 \quad \forall\, j
\end{aligned}
\label{eq:robust_lp}
\end{equation}

\subsection{Phase~3: The Fuzzy Programming Solution Approach}

This paper utilizes an efficient single-phase technique — the TH
method introduced by \citet{Torabi2008} — to deal with the concerned
MOLP model. Let the compact MOLP model be:

\begin{equation}
\begin{aligned}
  &\text{Min}\; Z_1,\quad \text{Min}\; Z_2,\quad \text{Max}\; Z_3 \\
  &\text{s.t.}\quad Ax \le b,\quad x \in X
\end{aligned}
\label{eq:molp}
\end{equation}

\textbf{Step~1:} Calculate the positive and negative ideal solutions
(PIS and NIS, respectively) for objective function $h$; $h = 1, 2,
3$.

\begin{equation}
Z_1^{\delta-\text{NIS}} = \text{Max}\{Z_1(x_2^{\delta-\text{PIS}}),\,
Z_1(x_3^{\delta-\text{PIS}})\}
\label{eq:nis1}
\end{equation}

\begin{equation}
Z_2^{\delta-\text{NIS}} = \text{Max}\{Z_2(x_1^{\delta-\text{PIS}}),\,
Z_2(x_3^{\delta-\text{PIS}})\}
\label{eq:nis2}
\end{equation}

\begin{equation}
Z_3^{\delta-\text{NIS}} = \text{Min}\{Z_3(x_1^{\delta-\text{PIS}}),\,
Z_3(x_2^{\delta-\text{PIS}})\}
\label{eq:nis3}
\end{equation}

\textbf{Step~2:} Calculate the satisfaction level of objective
function $h$; $h = 1, 2, 3$ under solution vector $x$:

\begin{equation}
\mu_1(x) = \begin{cases}
  1 & \text{if } Z_1 < Z_1^{\delta-\text{PIS}} \\[4pt]
  \dfrac{Z_1^{\delta-\text{NIS}} - Z_1}
        {Z_1^{\delta-\text{NIS}} - Z_1^{\delta-\text{PIS}}}
    & \text{if } Z_1^{\delta-\text{PIS}} \le Z_1 \le Z_1^{\delta-\text{NIS}}\\[6pt]
  0 & \text{if } Z_1 > Z_1^{\delta-\text{NIS}}
\end{cases}
\label{eq:mu1}
\end{equation}

\begin{equation}
\mu_3(x) = \begin{cases}
  1 & \text{if } Z_3 > Z_3^{\delta-\text{PIS}} \\[4pt]
  \dfrac{Z_3 - Z_3^{\delta-\text{NIS}}}
        {Z_3^{\delta-\text{PIS}} - Z_3^{\delta-\text{NIS}}}
    & \text{if } Z_3^{\delta-\text{PIS}} \le Z_3 \le Z_3^{\delta-\text{NIS}}\\[6pt]
  0 & \text{if } Z_3 < Z_3^{\delta-\text{NIS}}
\end{cases}
\label{eq:mu3}
\end{equation}

\textbf{Step~3:} Provide the single objective formulation based on
the TH aggregation function:

\begin{equation}
\begin{aligned}
  \text{Max}\;\lambda(x) &= \Gamma * \lambda_0 + (1-\Gamma)
    \sum_h \vartheta_h * \mu_h(x) \\
  \text{s.t.} \quad
    &\lambda_0 \le \mu_h(x),\quad h = 1,2,3 \\
    &x \in F(x)\lambda_0,\;\lambda \in [0,1] \\
    &\sum_h \vartheta_h = 1,\quad \vartheta_h \ge 0
\end{aligned}
\label{eq:th}
\end{equation}

% ============================================================
\section{Problem Statement}
% ============================================================

The designed [domain] supply chain network in this paper is based on
a real case study and it contains [N] echelons including [echelon 1],
[echelon 2], [echelon 3], [echelon 4], and [echelon 5], as shown
in Fig.~\ref{fig:network}.

\begin{itemize}
  \item \textbf{[Echelon 1]:} [Product] supply does not have any
        particular pattern and is tainted with uncertainty since
        [source] behavior is not completely known or predictable.

  \item \textbf{[Echelon 2]:} [Collection / production] process is
        generally covered by mobile and local facilities, which are
        distinguished by their capacities, functions, and costs.

  \item \textbf{[Echelon 3]:} Being received by regional centers,
        [product] units are tested to make sure they are safe against
        [hazards] and capable of being fractionated into sub-products.

  \item \textbf{[Echelon 4]:} The demand for [product] is accompanied
        with a high level of uncertainty and has no regular order.

  \item \textbf{[Echelon 5]:} The final destination of discarded
        [product] is a disposal center.
\end{itemize}

\subsection{Notations}

The following indices, parameters, and decision variables will be
used for the purpose of model formulations (uncertain parameters are
represented with a tilde on):

\subsubsection{Indices}

\begin{tabbing}
  \hspace{1.8cm}\= \kill
  $I$     \> Index of $i$th [donor/source] group ($i, i' \in I$).\\
  $J_1$   \> Index of locations for $j_1$th fix [facility]
             ($j_1 \in J_1$).\\
  $J_2$   \> Index of locations for $j_2$th mobile [facility]
             ($j_2 \in J_2$).\\
  $J$     \> Index of locations for $j$th [facility]
             ($j \in J,\, j = j_1 \cup j_2$).\\
  $K$     \> Index of locations for $k$th regional center
             ($k \in K$).\\
  $R$     \> Index of locations for $r$th disposal center
             ($r \in R$).\\
  $P$     \> Index of $p$th [product] ($p \in P$).\\
  $M$     \> Index of $m$th [method] ($m \in M$).\\
  $H$     \> Index of $h$th demand zone ($h \in H$).\\
  $T$     \> Index of $t$th time period ($t, t' \in T$).
\end{tabbing}

\subsubsection{Binary Variables}

\begin{tabbing}
  \hspace{1.8cm}\= \kill
  $Y_j$    \> 1 if $j$th [facility] $j$ is established; 0
              otherwise.\\
  $Z'_{jk}$\> 1 if $j$th [facility] is allocated to $k$th regional
              center; 0 otherwise.\\
  $Z_{ij}$ \> 1 if $i$th [source] group is allocated to $j$th
              [facility]; 0 otherwise.
\end{tabbing}

\subsection{Objective Functions}

The proposed multi-objective formulation is structured as follows:

\subsubsection{First Objective Function}
The first objective function seeks to minimize the total network
cost involving establishment cost (EC), operating cost (OC),
transportation costs (TC), inventory cost (IC), production cost
(PC) and unusable cost (UC) during multiple periods:

\begin{equation}
\text{Min}\; Z_1 = EC + OC + TC + IC + PC + UC
\label{eq:z1}
\end{equation}

\begin{equation}
EC = \sum_{j_1}\left(\tilde{E}_{j_1} + \tilde{F}_{j_1}
     + \tilde{LA}_{j_1}\right) * Y_{j_1}
     + \sum_{j_2}\left(\tilde{E}_{j_2} + \tilde{F}_{j_2}\right)
     * Y_{j_2}
\label{eq:ec}
\end{equation}

\begin{equation}
OC = \sum_{m}\sum_{i}\sum_{j}\sum_{t}
     \widetilde{OC}_m * V_{mijt}
\label{eq:oc}
\end{equation}

\begin{equation}
TC = \sum_{m}\sum_{j}\sum_{k}\sum_{t}
     \widetilde{TC}_{jk} * V'_{mjkt}
   + \sum_{k}\sum_{h}\sum_{p}\sum_{t}\sum_{t'}
     \widetilde{TC}'_{kh} * W_{khptt'}
\label{eq:tc}
\end{equation}

\begin{equation}
IC = \sum_{p}\sum_{k}\sum_{t}\sum_{t'}
     \widetilde{IC}_{pkt} * B_{pktt'}
   + \sum_{p}\sum_{h}\sum_{t}\sum_{t'}
     \widetilde{IC}'_{pht} * B'_{phtt'}
\label{eq:ic}
\end{equation}

\begin{equation}
PC = \sum_{p}\sum_{m}\sum_{k}\sum_{t}
     \widetilde{PC}_{pm} * Q_{pmkt}
\label{eq:pc}
\end{equation}

\begin{equation}
UC = \sum_{p}\sum_{k}\sum_{r}\sum_{t}
     \widetilde{TD}_{kr} * E_{pkrt}
   + \sum_{p}\sum_{k}\sum_{r}\sum_{t}
     \widetilde{TD}'_{kr} * E'_{pkrt'}
   + \sum_{p}\sum_{h}\sum_{r}\sum_{t}
     \widetilde{TD}''_{hr} * E''_{phrt'}
\label{eq:uc}
\end{equation}

\subsubsection{Second Objective Function}
The second objective function utilizes the following linear relation
of [product] freshness introduced by \citet{Osvald2008}:

\begin{equation}
\text{Min}\; Z_2 = \sum_p \varphi_p
\label{eq:z2}
\end{equation}

\begin{equation}
\varphi_p = \sum_k \sum_h \sum_t \sum_{t'}
  \vartheta_p * \text{Max}\!\left\{QLP_p,\,
  \frac{t'-t}{L_p}\right\}
  * \left(B_{pktt'} + B'_{phtt'}\right)
\label{eq:phi}
\end{equation}

where $\varphi_p$ is quality loss function of $p^{\text{th}}$
product, $QLP_p$ represents the point of quality reduction of
$p^{\text{th}}$ product, and $L_p$ defines product's shelf life.

\subsubsection{Third Objective Function}
Locating [facilities] is an intricate decision making involving
both qualitative and quantitative attributes:

\begin{equation}
\text{Max}\; Z_3 = \sum_{m}\sum_{i}\sum_{j_1}\sum_{t}
  QF_{j_1} * V_{m,i,j_1,t}
  + \sum_{m}\sum_{i}\sum_{j_2}\sum_{t}
  QF'_{j_2} * V_{m,i,j_2,t}
\label{eq:z3}
\end{equation}

\subsection{Model Constraints}

\begin{equation}
Z_{ij} \le Y_j \quad \forall\, i, j
\label{eq:c1}
\end{equation}

\begin{equation}
\sum_m \sum_i V_{mijt} \le \tilde{C}_j \quad \forall\, j,\, t
\label{eq:c2}
\end{equation}

\begin{equation}
\sum_m \sum_j V_{mijt} \le \tilde{S}_{it} \quad \forall\, i,\, t
\label{eq:c3}
\end{equation}

\begin{equation}
\delta_{ij} * Z_{ij} \le H \quad \forall\, i,\, j
\label{eq:c4}
\end{equation}

\begin{equation}
Z'_{jk}\!\left(\frac{\delta'_{jk}}{SP}\right) \le TI
  \quad \forall\, j,\, k
\label{eq:c5}
\end{equation}

\begin{equation}
\sum_i V_{mijt} = \sum_k V'_{mjkt} \quad \forall\, m,\, j,\, t
\label{eq:c6}
\end{equation}

\begin{equation}
\varphi_m * \delta_{pm} * \sum_j V'_{mjkt} = Q_{pmkt}
  \quad \forall\, p,\, m,\, k,\, t
\label{eq:c7}
\end{equation}

\begin{equation}
\sum_p PT_p * G_{pkt} \le \widetilde{OT}_{kt} \quad \forall\, k,\, t
\label{eq:c8}
\end{equation}

\begin{equation}
- \Psi \le X_{it} - X_{i't} \le \Psi \quad \forall\, i,\, i',\, t
\label{eq:c9}
\end{equation}

\begin{equation}
Z_{ij},\, Z'_{jk} \in \{0,1\} \quad \forall\, i,\, j,\, k
\label{eq:bin}
\end{equation}

\subsection{Robust Framework}

In this section, we utilize the robust optimization approach
described in Section~3.2 to account for the uncertain environment.
The robust counterpart of the deterministic programming model can be
formulated as follows:

\begin{equation}
\text{Min}\; \xi
\label{eq:robobj}
\end{equation}

\begin{multline}
\sum_{j_1}\left(E^E_{j_1} + F^F_{j_1} + LA^{LA}_{j_1}\right)
  + \sum_{j_2}\left(E^E_{j_2} + F^F_{j_2}\right)
  + \sum_m \sum_i \sum_j \sum_t P^{OC}_m * V_{mijt}\\
  + \sum_{m,j,k,t} P^{TC}_{j,k} * V'_{mjkt}
  + \cdots
  + \lambda^0 * \Gamma^0 \le \xi
\label{eq:robust_full}
\end{multline}

% ============================================================
\section{Case Description}
% ============================================================

[National / Regional Organization] has been working since [Year].
According to statistics, [Organization] includes [N1] [facilities]
and [N2] regional centers for testing and processing which are
located in [N3] provinces of the country \cite{Abolghasemi2009}.
To illustrate the usefulness of the proposed model, this paper
investigates a real case study of [City] which is the capital of
[Region] and the [Nth] most populated city of [Country] having a
population of about [X] million.

[City] consists of [N] districts, dispersed as shown in
Fig.~\ref{fig:map}, and plays a critical role in the provision of
[product] for the country. Based on the above descriptions, the
established [facilities] are in charge of collecting [product] from
the [source] points. The properties of each district are presented
in Table~\ref{tab:districts}. The properties of each candidate
[facility] and each hospital are presented in
Tables~\ref{tab:facilities} and~\ref{tab:hospitals}, respectively.

With respect to [relevant standard], [product] demand can be
estimated based on the patients' population or according to each
hospital's bed. Transportation cost between each pair of facilities
is proportional to the traveling distance, estimated to be [X]\$ per
kilometer.

% ============================================================
\section{Results and Discussions}
% ============================================================

In this section, the first phase is implemented by applying software
[Name] for the purpose of obtaining the importance degree of each
candidate local and mobile location.

\textbf{Step~1. Calculating the weights of the DMs:} in this part,
the importance degree of the DMs is computed according to
Section~3.1.2. Using Eqs.~\eqref{eq:ek} and~\eqref{eq:gamma}, the
weights of the [N] DMs are obtained as follows:
$\Gamma_1 = 0.161$, $\Gamma_2 = 0.173$, $\Gamma_3 = 0.138$,
$\Gamma_4 = 0.166$, $\Gamma_5 = 0.194$, and $\Gamma_6 = 0.168$.

\textbf{Step~2. Calculating the weights of the attributes:} this
part obtains the importance degree of each attribute based on
model~\eqref{eq:maxq}. According to the DMs' preference-intervals
for each attribute, say $I = \{0.1 \le \omega_1 \le 0.2,\;
0.1 \le \omega_2 \le 0.3, \ldots\}$, the vector of attributes
weights are computed as: $\omega = \{0.2, 0.1, 0.2, 0.3, 0.1, 0.1\}$.

\textbf{Step~3. Identifying the candidate locations:} in this part,
the most qualified candidate locations are identified and ranked
based on the obtained weights.

In continue, the computational experiments are carried out through
CPLEX solver by utilizing software GAMS on a PC with Intel Core i7,
CPU 2.4~GHz and 16~GB of RAM. As shown in Table~\ref{tab:results},
the results obtained by robust and deterministic models are
presented under the variability of uncertain parameters and
conservatism level. For this purpose, we first set the values of
the compensation coefficient ($\Upsilon$) to 0.4 and the relative
importance of the objective functions ($\vartheta_1, \vartheta_2,
\vartheta_3$) to 0.6, 0.2, and 0.2, respectively.

A general expectation is that the network cost of the robust
approach under parameters variability is higher than that of the
deterministic mode since it imposes additional costs of handling
uncertainty and making a stable network. By increasing the level of
uncertainty, the cost increase becomes even more tangible.
Furthermore, the influence of utilizing robust approach on computing
time is negligible comparing to the deterministic attitude.

\subsection{The Supply Chain Design Under Different Variability
Levels of Uncertain Parameters}

The impact of uncertainty levels on optimal structure of the supply
chain provided by robust and deterministic approaches is mapped and
compared in Fig.~\ref{fig:sensitivity}. A general understanding is
that in addition to the existing local [facilities], by increasing
the variability from 10\% to 50\%, a slight increase in the number
of mobile [facilities] is justified to cover the realized
perturbation.

\subsection{The Model Performance Under Different Levels of
Compensation Coefficients ($\Gamma$)}

In this part, the impact of varying $\Gamma$ on the performance of
the model is evaluated. As the value of $\Gamma$ increases the model
forces the membership functions to get closer to each other while on
the other hand, the lower levels of $\Gamma$ give the importance to
the objective with higher weight.

\subsection{The Model Performance Under Different Levels of
Conservatism and Data Variability}

This section aims to analyze the behavior of the objective function
of cost while adopting different risk-averse attitudes and variability
levels. The analysis is presented in three parts: (a) Sensitivity of
the model to an uncertain production cost, (b) Sensitivity of the
model to an uncertain capacity of a [facility], and (c) Sensitivity
of the model to uncertain demand.

% ============================================================
\section{Discussions and Conclusions}
% ============================================================

Designing [domain] supply chain networks is tainted with a great
degree of intricacies due to properties including rather irregular
supply and demand, wastage and shortage costs as well as
[domain]-related considerations (e.g., perishability of [product]),
which are not easily tractable. With this significance in mind, this
paper proposes a novel 3P-MAMOR model to formulate the importance of
the factors above to best match the real world conditions. The
proposed model aims to consider both cost-oriented and
quality-oriented approaches by minimizing the overall network cost,
preserving as much as possible the [product freshness], and seeking
the most qualified locations for [facilities]. This paper also
applies a real case study to accommodate the real world application
of the proposed 3P-MAMOR model.

According to the results obtained from the solution, it seems
imperative to make a tradeoff between cost-oriented and
quality-oriented approaches. Qualitative factors like availability,
market environment, and transportation conditions have significant
effect on finding the preferable locations for collection facilities.
Moreover, the robust approach has dominated the deterministic model
in terms of the mean and SD of the value of the total network cost.

Future research can be directed in investigating the problem while
using other MADM techniques such as fuzzy AHP or fuzzy TOPSIS to
evaluate the critical factors involved. The interested practitioners
can also develop the hospital level by accounting for compatibility
procedures. Moreover, the mathematical modelling in realm of
[donation system] could be developed by considering age-based
variety of [products].

% ============================================================
%  TABLES
% ============================================================

% ---- Table 1: Nomenclature (2-column wide) -----------------
\begin{table*}[ht]
\caption{Nomenclature of the investigated papers' characteristics.}
\label{tab:nomenclature}
\tabfont
\setlength{\tabcolsep}{5pt}
\begin{tabular}{p{2.5cm}lp{2.5cm}lp{2.5cm}l}
\toprule
\multicolumn{2}{l}{\headfont Hierarchical level}
 & \multicolumn{2}{l}{\headfont Modelling approach}
 & \multicolumn{2}{l}{\headfont Performance measure} \\
\midrule
Collection    & CO & Simulation             & S   & Cost       & C  \\
Production    & PR & Statistical analysis   & SA  & Distance   & D  \\
Inventory     & IN & (Stochastic) DP        & (S)DP & Shortage  & SH \\
Distribution  & DI & Integer Programming    & IP  & Outdates   & O  \\
Integrated    & INT& (M)INLP                & (M)INLP & Freshness & F \\
\midrule
\multicolumn{2}{l}{\headfont Planning horizon}
 & \multicolumn{2}{l}{\headfont Type of blood product}
 & \multicolumn{2}{l}{\headfont Uncertainty approach} \\
\midrule
Strategic     & STR & Whole blood            & WB  & Fuzzy      & FZ \\
Tactical      & TAC & Platelet               & PLT & Stochastic & ST \\
Operational   & OPR & Plasma                 & PLS & Robust     & RO \\
\midrule
\multicolumn{2}{l}{\headfont Solution method}
 & \multicolumn{2}{l}{\headfont Objective function}
 & \multicolumn{2}{l}{} \\
\midrule
Exact         & EX & Single-objective       & SO  &            & \\
(Meta)heuristic & ME & Multi-objective       & MO  &            & \\
\bottomrule
\end{tabular}
\end{table*}

% ---- Table 2: District properties --------------------------
\begin{table}[H]
\caption{The properties of each district.}
\label{tab:districts}
\tabfont
\setlength{\tabcolsep}{4pt}
\begin{tabular}{lrrrp{2.0cm}}
\toprule
\headfont District & \headfont Population
  & \headfont Rate (\%) & \headfont Supply
  & \headfont (LONG$_i$, LAT$_i$) \\
\midrule
1  & 212,206 & 14 & 38 & (38.077, 46.311) \\
2  & 169,047 & 11 & 29 & (38.063, 46.287) \\
3  & 243,400 & 16 & 42 & (38.067, 46.255) \\
4  & 316,126 & 21 & 55 & (38.076, 46.253) \\
5  &  92,274 &  6 & 16 & (38.048, 46.366) \\
6  &  94,897 &  6 & 16 & (38.086, 46.232) \\
7  & 143,460 & 10 & 27 & (38.043, 46.191) \\
8  &  28,700 &  2 &  6 & (38.081, 46.298) \\
9  &     324 &  0 &  1 & (38.014, 46.423) \\
10 & 194,564 & 13 & 34 & (38.112, 46.290) \\
\bottomrule
\end{tabular}
\end{table}

% ---- Table 3: Candidate facility properties ----------------
\begin{table}[H]
\caption{The properties of each candidate local and mobile
[facility].}
\label{tab:facilities}
\tabfont
\setlength{\tabcolsep}{4pt}
\begin{tabular}{lp{3.2cm}lp{2.0cm}}
\toprule
\headfont Row & \headfont Facility
  & \headfont Type & \headfont (LONG$_i$, LAT$_i$) \\
\midrule
1 & [Facility Name 1] & Candidate local  & (38.073, 46.252) \\
2 & [Facility Name 2] & Candidate local  & (38.073, 46.303) \\
3 & [Facility Name 3] & Candidate local  & (38.066, 46.311) \\
4 & [Facility Name 4] & Candidate local  & (38.076, 46.288) \\
5 & [Facility Name 5] & Candidate local  & (38.065, 46.271) \\
6 & [Facility Name 6] & Candidate local  & (38.059, 46.348) \\
\midrule
1 & [Mobile 1]        & Candidate mobile & (38.055, 46.278) \\
2 & [Mobile 2]        & Candidate mobile & (38.061, 46.333) \\
3 & [Mobile 3]        & Candidate mobile & (38.084, 46.265) \\
4 & [Mobile 4]        & Candidate mobile & (38.063, 46.303) \\
5 & [Mobile 5]        & Candidate mobile & (38.091, 46.264) \\
\bottomrule
\end{tabular}
\end{table}

% ---- Table 4: Demand zone properties -----------------------
\begin{table}[H]
\caption{The properties of each demand zone.}
\label{tab:hospitals}
\tabfont
\setlength{\tabcolsep}{4pt}
\begin{tabular}{lp{2.8cm}rp{2.0cm}}
\toprule
\headfont Row & \headfont Hospital
  & \headfont Beds & \headfont (LONG$_i$, LAT$_i$) \\
\midrule
1  & [Hospital 1]  & 599 & (38.108, 46.302) \\
2  & [Hospital 2]  &  67 & (38.093, 46.279) \\
3  & [Hospital 3]  &  85 & (38.076, 46.259) \\
4  & [Hospital 4]  & 154 & (38.077, 46.279) \\
5  & [Hospital 5]  & 307 & (38.082, 46.306) \\
6  & [Hospital 6]  & 168 & (38.080, 46.306) \\
7  & [Hospital 7]  &  56 & (38.069, 46.297) \\
8  & [Hospital 8]  &  37 & (38.068, 46.297) \\
9  & [Hospital 9]  & 124 & (38.060, 46.282) \\
10 & [Hospital 10] &  93 & (38.078, 46.318) \\
\bottomrule
\end{tabular}
\end{table}

% ---- Table 5: Main results table (2-col wide) --------------
\begin{table*}[ht]
\caption{Summary of results under uncertain parameters and
deterministic condition ($\Gamma=0.4$, $\vartheta_1=0.6$,
$\vartheta_2=\vartheta_3=0.2$).}
\label{tab:results}
\tabfont
\setlength{\tabcolsep}{4pt}
\renewcommand{\arraystretch}{1.20}
\begin{tabular}{lccccccc}
\toprule
\multirow{2}{*}{\headfont Var.}
 & \headfont Deterministic $Z_1$
 & \multicolumn{4}{c}{\headfont Robust $Z_1$ --- Conservatism level}
 & \headfont Det.\ $Z_2$
 & \headfont Det.\ time (s) \\
\cmidrule(lr){3-6}
 & & \headfont 60\% & \headfont 70\% & \headfont 80\% & \headfont 90\%
 & & \\
\midrule
0\%  & 682,918 & \multicolumn{4}{c}{---}      & 688 & 72  \\
10\% & 698,008 & 683,849 & 690,235 & 671,599 & 661,632 & 589 & --- \\
20\% & 710,753 & 709,518 & 697,482 & 678,357 & 658,549 & 491 & 190 \\
30\% & 708,092 & 695,032 & 739,653 & 708,092 & 728,294 & 305 & 282 \\
40\% & 761,267 & 744,372 & 742,249 & 663,780 & 767,346 & 586 & 231 \\
50\% & 775,397 & 794,836 & 786,078 & 680,403 & 773,581 & 345 & 206 \\
\bottomrule
\end{tabular}
\end{table*}

% ---- Table 6: Compensation coefficients --------------------
\begin{table}[H]
\caption{The model performance under different levels of
compensation coefficients~($\Gamma$).}
\label{tab:gamma}
\tabfont
\setlength{\tabcolsep}{3pt}
\begin{tabular}{crrrrrrrr}
\toprule
\headfont Prob. & $\Gamma$ & $\vartheta_1,\vartheta_2,\vartheta_3$
  & $Z_1$ & $Z_2$ & $Z_3$
  & $\mu_1$ & $\mu_2$ & $\mu_3$ \\
\midrule
1 & 0.0 &              & 331,717 & 410 & 511 & 0.995 & 0.996 & 0.355 \\
2 & 0.1 &              & 331,698 & 413 & 511 & 0.995 & 0.996 & 0.355 \\
3 & 0.2 &              & 421,223 & 418 & 592 & 0.944 & 0.996 & 0.468 \\
4 & 0.3 &              & 708,844 & 348 & 812 & 0.781 & 0.996 & 0.776 \\
5 & 0.4 & $\vartheta_1=0.7$& 709,501 & 343 & 812 & 0.780 & 0.996 & 0.776 \\
6 & 0.5 & $\vartheta_2=0.2$ & 706,425 & 418 & 812 & 0.780 & 0.996 & 0.776 \\
7 & 0.6 & $\vartheta_3=0.1$ & 706,425 & 418 & 812 & 0.781 & 0.996 & 0.776 \\
8 & 0.7 &              & 706,425 & 418 & 812 & 0.781 & 0.996 & 0.776 \\
9 & 0.8 &              & 706,425 & 418 & 812 & 0.781 & 0.996 & 0.776 \\
10& 0.9 &              & 706,322 & 440 & 812 & 0.781 & 0.995 & 0.776 \\
11& 1.0 &              & 716,357 & 11884 & 812 & 0.776 & 0.874 & 0.776 \\
\bottomrule
\end{tabular}
\end{table}

% ============================================================
%  FIGURES
% ============================================================

% ---- Figure 1: 3P-MAMOR framework flowchart ----------------
\begin{figure*}[ht]
\centering
\begin{tikzpicture}[
  box/.style={rectangle, draw=black, rounded corners=2pt,
              minimum width=5.5cm, minimum height=0.55cm,
              font=\footnotesize, align=center, fill=white},
  arr/.style={-Stealth, thick},
  phase/.style={rectangle, draw=black, dashed, rounded corners=3pt,
                inner sep=6pt, font=\footnotesize\bfseries}
]
% ---- Phase 1 boxes -----------------------------------------
\node[box] (id1) at (0,0)
  {Determine the decision matrix for candidate [facilities]};
\node[box] (id2) at (-2.5,-1.0)
  {Determine the weights of DMs};
\node[box] (id3) at (2.5,-1.0)
  {Determine the weights of attributes};
\node[box] (id4) at (0,-2.0)
  {Rank the candidate [facilities]};
\node[box] (id5) at (0,-3.0)
  {Identify the importance of candidate [facilities]};

\draw[arr] (id1) -- (id2);
\draw[arr] (id1) -- (id3);
\draw[arr] (id2) -- (id4);
\draw[arr] (id3) -- (id4);
\draw[arr] (id4) -- (id5);

\node[phase, fit=(id1)(id2)(id3)(id4)(id5),
      label={[font=\footnotesize\bfseries]left:\textbf{PHASE 1 (Identification)}}]
  (ph1) {};

% ---- Phase 2 boxes -----------------------------------------
\node[box] (pr1) at (0,-4.2)
  {Describe the multi-objective deterministic model};
\node[box] (pr2) at (0,-5.0)
  {Define the uncertain parameters};
\node[box] (pr3) at (0,-5.8)
  {Formulate the multi-objective robust model};

\draw[arr] (id5) -- (pr1);
\draw[arr] (pr1) -- (pr2);
\draw[arr] (pr2) -- (pr3);

\node[phase, fit=(pr1)(pr2)(pr3),
      label={[font=\footnotesize\bfseries]left:\textbf{PHASE 2 (Protection)}}]
  (ph2) {};

% ---- Phase 3 boxes -----------------------------------------
\node[box] (in1) at (0,-7.0)
  {Determine positive and negative ideal solutions};
\node[box] (in2) at (0,-7.8)
  {Calculate the membership function for each objective};
\node[box] (in3) at (0,-8.6)
  {Convert the multi-objective model to single-objective};
\node[box] (in4) at (0,-9.4)
  {Obtain the efficient solution for each objective};

\draw[arr] (pr3) -- (in1);
\draw[arr] (in1) -- (in2);
\draw[arr] (in2) -- (in3);
\draw[arr] (in3) -- (in4);

\node[phase, fit=(in1)(in2)(in3)(in4),
      label={[font=\footnotesize\bfseries]left:\textbf{PHASE 3 (Integration)}}]
  (ph3) {};

\end{tikzpicture}
\caption{Visual representation of the 3P-MAMOR modelling approach.}
\label{fig:framework}
\end{figure*}

% ---- Figure 2: Supply chain network schematic --------------
\begin{figure*}[ht]
\centering
\begin{tikzpicture}[
  node distance=1.5cm,
  entity/.style={rectangle, draw, rounded corners=2pt,
                 minimum width=2.0cm, minimum height=0.6cm,
                 font=\scriptsize, align=center, fill=white},
  phase_box/.style={rectangle, draw=gray, dashed, rounded corners=3pt,
                    font=\scriptsize\bfseries, inner sep=4pt,
                    fill=gray!5},
  arr/.style={-Stealth, thick, draw=black}
]
% Donors
\node[entity] (donors) at (0,0) {Donors ($I$)};

% Collection facilities
\node[entity] (fix)    at (2.8, 0.6) {Fix Collection\\ Centers ($J_1$)};
\node[entity] (mobile) at (2.8,-0.6) {Mobile Collection\\ Centers ($J_2$)};

% Regional blood center
\node[entity] (rbc) at (6.0, 0) {Regional Blood\\ Centers ($K$)};

% Demand zones
\node[entity] (demand) at (9.2, 0.6) {Demand Zones\\ ($H$)};

% Disposal
\node[entity] (disposal) at (9.2,-0.6) {Disposal\\ Centers ($R$)};

% Arrows
\draw[arr] (donors) -- (fix);
\draw[arr] (donors) -- (mobile);
\draw[arr] (fix)    -- (rbc);
\draw[arr] (mobile) -- (rbc);
\draw[arr] (rbc)    -- (demand);
\draw[arr] (rbc)    -- (disposal);
\draw[arr] (demand.south) -- ++(0,-0.4) -| (disposal.north);

% Phase annotations
\node[phase_box] at (1.4,-1.6) {\textbf{Collection Phase}};
\node[phase_box] at (6.0,-1.6)
  {\textbf{Production and Inventory Phases}};
\node[phase_box] at (9.2,-1.6) {\textbf{Distribution Phase}};

% Inset labels
\node[font=\scriptsize, align=center] at (6.0, 0.85)
  {\textit{Laboratory — Process — Storage}};

\end{tikzpicture}
\caption{The concerned [domain] supply chain overview.}
\label{fig:network}
\end{figure*}

% ---- Figure 3: Sensitivity line chart with error bars ------
\begin{figure}[H]
\centering
\begin{tikzpicture}
\begin{axis}[
  width         = \columnwidth,
  height        = 5.2cm,
  xlabel        = {Variability (\%)},
  ylabel        = {Total supply chain costs (\%)},
  xtick         = {10,20,30,40,50},
  xticklabels   = {10\%,20\%,30\%,40\%,50\%},
  ymin          = 600000, ymax = 1250000,
  ytick         = {600000,680000,760000,840000,920000,
                   1000000,1080000,1160000,1240000},
  yticklabel style = {font=\scriptsize,
    /pgf/number format/fixed,
    /pgf/number format/precision=0,
    /pgf/number format/fixed zerofill=false},
  xlabel style  = {font=\footnotesize},
  ylabel style  = {font=\footnotesize},
  tick label style = {font=\scriptsize},
  legend style  = {font=\scriptsize, at={(0.98,0.05)},
                   anchor=south east, draw=none},
  axis lines    = left,
  grid          = none,
  line width    = 1pt,
]
% Deterministic — blue line with error bars (mean + SD)
\addplot[color=elsblue, mark=*, mark size=2.5pt,
         error bars/.cd, y dir=both, y explicit]
  coordinates {
    (10,  758190) +- (0,27503)
    (20,  881051) +- (0,80261)
    (30,  983367) +- (0,105669)
    (40, 1037946) +- (0,128225)
    (50, 1037946) +- (0,128225)
  };
\addlegendentry{Deterministic}

% Robust — red/dark line with error bars
\addplot[color=red!70!black, mark=square*, mark size=2.0pt,
         error bars/.cd, y dir=both, y explicit]
  coordinates {
    (10, 661632) +- (0,12125)
    (20, 658549) +- (0,35477)
    (30, 685812) +- (0,51530)
    (40, 665128) +- (0,46375)
    (50, 647777) +- (0,56237)
  };
\addlegendentry{Robust}
\end{axis}
\end{tikzpicture}
\caption{The mean and standard deviation of the first objective
function value (Conservatism Level 90\%).}
\label{fig:sensitivity}
\end{figure}

% ---- Figure 4: Schematic map of city districts -------------
\begin{figure}[H]
\centering
\begin{tikzpicture}[scale=0.70]
% Simplified district shapes as irregular polygons
% mimicking a geographic district map
\fill[gray!20] (0,0) -- (3,0) -- (3.5,1.5) -- (2,2.5)
               -- (0.5,2) -- cycle;
\fill[gray!30] (3,0) -- (5.5,0) -- (5.5,2) -- (3.5,1.5) -- cycle;
\fill[gray!15] (0.5,2) -- (2,2.5) -- (2,4) -- (0,4) -- cycle;
\fill[gray!25] (2,2.5) -- (3.5,1.5) -- (4,3) -- (2,4) -- cycle;
\fill[gray!10] (3.5,1.5) -- (5.5,2) -- (5.5,4) -- (4,3) -- cycle;
\fill[gray!35] (0,4) -- (2,4) -- (2,5.5) -- (0,5.5) -- cycle;
\fill[gray!20] (2,4) -- (4,3) -- (4.5,5.5) -- (2,5.5) -- cycle;
\fill[gray!15] (4,3) -- (5.5,4) -- (5.5,5.5) -- (4.5,5.5) -- cycle;

% District borders
\draw[black, line width=0.5pt]
  (0,0) -- (3,0) -- (3.5,1.5) -- (2,2.5) -- (0.5,2) -- cycle
  (3,0) -- (5.5,0) -- (5.5,2) -- (3.5,1.5) -- cycle
  (0.5,2) -- (2,2.5) -- (2,4) -- (0,4) -- cycle
  (2,2.5) -- (3.5,1.5) -- (4,3) -- (2,4) -- cycle
  (3.5,1.5) -- (5.5,2) -- (5.5,4) -- (4,3) -- cycle
  (0,4) -- (2,4) -- (2,5.5) -- (0,5.5) -- cycle
  (2,4) -- (4,3) -- (4.5,5.5) -- (2,5.5) -- cycle
  (4,3) -- (5.5,4) -- (5.5,5.5) -- (4.5,5.5) -- cycle;

% District labels
\node[font=\tiny] at (1.5,0.8)  {\textbf{3}};
\node[font=\tiny] at (4.2,0.8)  {\textbf{4}};
\node[font=\tiny] at (0.8,3.0)  {\textbf{7}};
\node[font=\tiny] at (2.8,2.8)  {\textbf{2}};
\node[font=\tiny] at (4.5,2.5)  {\textbf{5}};
\node[font=\tiny] at (1.0,4.8)  {\textbf{6}};
\node[font=\tiny] at (3.0,4.8)  {\textbf{1}};
\node[font=\tiny] at (4.8,4.8)  {\textbf{8}};

% Facility markers (circles = regional, squares = local, triangles = mobile)
\node[circle,fill=black,inner sep=2.5pt,
      label={[font=\tiny]right:RBC}] at (2.5,2.0) {};
\node[rectangle,fill=black,minimum size=5pt,
      label={[font=\tiny]right:LBC}] at (1.2,1.2) {};
\node[regular polygon,regular polygon sides=3,fill=black,
      minimum size=5pt,
      label={[font=\tiny]right:MBC}] at (3.8,1.8) {};

% Legend
\node[font=\tiny] at (2.8,-0.6)
  {\tikz\node[circle,fill=black,inner sep=2pt]{};~Regional Blood Center
   \quad
   \tikz\node[rectangle,fill=black,minimum size=4pt]{};~Local Blood Center
   \quad
   \tikz\node[regular polygon,regular polygon sides=3,
               fill=black,minimum size=4pt]{};~Mobile Blood Center};
\end{tikzpicture}
\caption{Geographic dispersion of districts of [City].}
\label{fig:map}
\end{figure}

% ============================================================
%  AUTHOR CONTRIBUTIONS & SUPPLEMENTARY
% ============================================================
\vspace{4pt}
\noindent\textbf{Appendix A. Supplementary material}

\noindent Supplementary data to this article can be found online at
\url{https://doi.org/10.1016/j.cie.2019.XX.XXX}.

% ============================================================
%  REFERENCES — Elsevier numbered style (Vancouver)
% ============================================================
\balance
\vspace{4pt}

\begin{thebibliography}{99}
\setlength{\itemsep}{1pt}
\setlength{\parskip}{0pt}

\bibitem{Author2015}
Author, A., \& Author, B. (2015).
Title of the referenced article.
\textit{Journal Name}, \textit{78}, 259--270.
\url{https://doi.org/10.1016/XXXX}

\bibitem{WHO2016}
WHO. (2016).
Blood safety and availability.
\textit{World Health Organization Report}.

\bibitem{Tofighi2016}
Tofighi, S., Ali~Torabi, S., \& Afshin~Mansouri, S. (2016).
Humanitarian logistics network design under mixed uncertainty.
\textit{European Journal of Operational Research},
\textit{250}(1), 239--250.
\url{https://doi.org/10.1016/j.ejor.2015.08.049}

\bibitem{Torabi2015}
Torabi, S.~A., Baghersad, M., \& Mansouri, S.~A. (2015).
Resilient supplier selection and order allocation under operational
and disruption risks.
\textit{Transportation Research Part E},
\textit{79}, 22--48.

\bibitem{Samani2018}
Samani, M.~R.~G., Torabi, S.~A., \& Hosseini-Motlagh, S.-M.\ (2018).
Integrated blood supply chain planning for disaster relief.
\textit{International Journal of Disaster Risk Reduction},
\textit{27}, 168--188.

\bibitem{Torabi2008}
Torabi, S.~A., \& Hassini, E. (2008).
An interactive possibilistic programming approach for multiple
objective supply chain master planning.
\textit{Fuzzy Sets and Systems},
\textit{159}(2), 193--214.

\bibitem{Belien2012}
Beli\"{e}n, J., \& Forc\'{e}, H. (2012).
Supply chain management of blood products: A literature review.
\textit{European Journal of Operational Research},
\textit{217}(1), 1--16.

\bibitem{Osorio2015}
Osorio, A.~F., Brailsford, S.~C., \& Smith, H.~K. (2015).
A structured review of quantitative models in the blood supply chain.
\textit{International Journal of Production Research},
\textit{53}(24), 7191--7212.

\bibitem{Bertsimas2003}
Bertsimas, D., \& Sim, M. (2003).
Robust discrete optimization and network flows.
\textit{Mathematical Programming},
\textit{98}(1--3), 49--71.

\bibitem{Bertsimas2004}
Bertsimas, D., \& Sim, M. (2004).
The price of robustness.
\textit{Operations Research},
\textit{52}(1), 35--53.

\bibitem{Ye2013}
Ye, J. (2013).
Multiple attribute group decision-making methods with unknown
weights in intuitionistic fuzzy setting and interval-valued
intuitionistic fuzzy setting.
\textit{International Journal of General Systems},
\textit{42}(5), 489--502.

\bibitem{Osvald2008}
Osvald, A., \& Stirn, L.~Z. (2008).
A vehicle routing algorithm for the distribution of fresh vegetables
and similar perishable food.
\textit{Journal of Food Engineering},
\textit{85}(2), 285--295.

\bibitem{Abolghasemi2009}
Abolghasemi, H., Maghsudlu, M., Amini~Kafi-Abad, S.,
\& Cheraghali, A. (2009).
Introduction to Iranian blood transfusion organization and blood
safety in Iran.
\textit{Iranian Journal of Public Health},
\textit{38}(1), 82--87.

\bibitem{Author2017}
Author, R., \& Author, Z. (2017).
Blood supply chain network design considering blood group
compatibility under uncertainty.
\textit{International Journal of Production Research},
\textit{55}(7), 2013--2033.

\bibitem{Author2014b}
Author, A., Author, B., \& Author, C. (2014).
Dynamic supply chain network design for the supply of blood in
disasters: A robust model with real world application.
\textit{Transportation Research Part E},
\textit{70}, 225--244.

\end{thebibliography}

\end{document}
